\section{A trace quadratic for an exceptional end}
Throughout this section the field has odd characteristic. Traces refer to
specified determinant-one lifts. Thus a projective element of order three
has chosen trace $1$, an element of order four has chosen trace
$\sqrt2$, and the two classes of order-five elements have traces
$(-1\pm\sqrt5)/2$. The last convention is used away from characteristic
five. Changing the sign of a lift changes none of the projective statements.

\begin{theorem}\label{E.TraceExtension}
Let $F$ be a finite field of odd order, and let $a,b\in F^\times$, $c\in F$ with
$d=a^2+b^2-4\ne0$. Up to simultaneous projective conjugacy, pairs
$A,B\in\mathrm{SL}_2(F)$ with traces $a,b$ and $\operatorname{tr}(AB)=0$
form one class. For each such pair, determinant-one matrices $C$ with
\[
 \operatorname{tr}(C)=c,\qquad
 \operatorname{tr}(BC)=\operatorname{tr}(ABC)=0
\]
are in bijection with the square roots in $F$ of
\[
 \Delta=a^2c^2-4(a^2+b^2+c^2-4).
\]
There are two, one or zero extensions according as $\Delta$ is a nonzero
square, zero or a nonsquare. Their invariant is
$2\operatorname{tr}(AC)-ac=\pm\sqrt\Delta$.
If $a,c\ne0$ and $c^2\ne4$, a projective conjugation realizes the
rotation reversal precisely when $\Delta=0$.
\end{theorem}
\begin{proof}
Put $J=\left(\begin{smallmatrix}0&1\\-1&0\end{smallmatrix}\right)$.
Choose $U,V\in F$ with $U^2+V^2=d$. Such a choice exists because the
sets of squares and of $d$ minus squares each have $(|F|+1)/2$ elements.
Set $u=(U+a)/2$, $w=(V-b)/2$ and
\[
 A=\begin{pmatrix}u&w+b\\w&a-u\end{pmatrix},\qquad B=A^{-1}J.
\]
Direct expansion gives determinant one, the required traces, and $AB=J$.
The nonzero $d$ says that $A,B$ have no common eigenline: equivalently,
$\det(AB-BA)=-d$. More generally the four matrices $I,A,B,AB$ are
linearly independent. This follows by diagonalizing a nonscalar $A$:
both off-diagonal entries of $B$ are nonzero, so $I,A$ span the diagonal
matrices and $B,AB$ supply the off-diagonal ones. In the remaining Jordan
case its unique eigenline is not $B$-invariant, and the same four-entry
calculation proves independence.
The Cayley--Hamilton identities and
$AB+BA=bA+aB-abI$ determine their multiplication table from $a,b$.
Consequently any second pair with the same three traces defines an
algebra automorphism of $M_2(F)$ by matching these bases. Such an
automorphism is inner: take a nonzero vector in the image of $E_{11}$,
and its image under the transformed $E_{21}$, as a basis. The transformed
matrix units then act as the usual matrix units. This proves pair uniqueness.

Write $C=\left(\begin{smallmatrix}r&s\\s&c-r\end{smallmatrix}\right)$,
as is forced by $\operatorname{tr}(JC)=0$, and put $R=2r-c$, $S=2s$.
The other trace and determinant equations are exactly
\[
 VR-US=-bc,\qquad R^2+S^2=c^2-4.
\]
Since $d\ne0$, every solution is uniquely of the form
\[
 R=\frac{-bcV+zU}{d},\qquad S=\frac{bcU+zV}{d},\qquad z^2=\Delta.
\]
This also gives $2\operatorname{tr}(AC)-ac=UR+VS=z$.
The projective centralizer of the ordered pair is trivial, so the
extensions have not acquired an additional conjugacy identification.

For the last assertion, nonzero traces remove the signs in determinant-one
lifts of $XAX^{-1}=A^{-1}$ and $XCX^{-1}=C$.
Put $A_0=A-aI/2$, $C_0=C-cI/2$. The centralizer algebra of $C$ is
$F[I,C_0]$, so $X=xI+yC_0$. The first equation is
\[
 XA_0+A_0X=2xA_0+y\operatorname{tr}(A_0C_0)I=0.
\]
Since $A_0$ is nonscalar, invertibility forces $x=0$, $y\ne0$ and
$\operatorname{tr}(A_0C_0)=z/2=0$. Conversely, if $z=0$, choose $X=C_0$.
It is invertible because $c^2\ne4$, has scalar square, commutes with $C$,
and inverts $A$ projectively. Also $C_0J=-JC_0$ by the trace condition,
so it fixes $AB=J$ projectively and sends $B$ to $A^2B$.
This is exactly the rotation reversal rule.
\end{proof}

\begin{lemma}\label{E.OddDegreeExclusion}
If $q=p^d\equiv3\pmod4$ and $d>1$ is a prime power, then $L(q)$
cannot be a chiral rotation group.
\end{lemma}
\begin{proof}
Only rank four remains. Neither end group fixes a point: it contains the
involution given by the product of its rotations, whereas an involution
of $L(q)$ has no rational fixed point when $q\equiv3\pmod4$.
By the retained natural-subfield maximal-subgroup classification, each
proper end lies in a conjugate of $L(p^{e_i})$ with $e_i\mid d$ proper.
Here the exceptional $A_4,S_4,A_5$ alternatives also lie in proper
subfield overgroups; the odd degree excludes a natural $P(p^{e_i})$.
Conjugate the first end to its standard position and choose lifts $A,B$
there. A lift of the third rotation has trace $c\in F_{p^{e_2}}$, since
trace is conjugacy invariant. Put $F=F_{p^{\operatorname{lcm}(e_1,e_2)}}$;
this is a proper subfield because the divisors of the prime power $d$
are linearly ordered.

The first pair is absolutely irreducible. Otherwise its common
algebraic eigenline is rational, contrary to the preceding fixed-point
observation, or its Frobenius conjugate is a distinct common eigenline.
In the latter case the pair is diagonal on the two lines and commutes.
Then $(AB)^2=1$ makes $A^2=B^{-2}$ projectively; the cyclic intersection
axiom forces both rotations to have order at most two, already excluded.
Thus $a^2+b^2-4\ne0$. Normalize $AB$ to $J$ over $F$ and apply the
explicit extension formulas. All entries of $C$ lie in $F(\sqrt\Delta)$.
The ambient extension has odd degree, so this quadratic extension is
$F$ itself. Hence all three rotations lie in the proper $L(F)$,
a contradiction. This replaces the source's quadratic calculation,
including its printed coefficient errors, by a nonsingular formula.
\end{proof}
